\documentclass[11pt]{article}
\usepackage{amscd}
\usepackage{amsfonts}
\usepackage{amsmath}
\usepackage{amssymb}
\usepackage{amsthm}
\usepackage{array}
\usepackage{color}
\usepackage{graphicx}
\usepackage{latexsym}
\usepackage{multicol}
\usepackage{multirow}
\usepackage{rotating}
\usepackage{subfigure}
\usepackage{ulem}
\usepackage{verbatim}
%RH LaTex Macros Apr 2014
%% query in the margin macro
\newcommand{\query}[1]{{$\bullet$}\marginpar{%
\vskip-\baselineskip %raise the marginpar a bit
\raggedright\footnotesize
%\itshape\hrule\smallskip#1\par\smallskip\hrule}}
\itshape\hrule\smallskip\parbox{.75in}{{$\bullet$} #1}\par\smallskip\hrule}}
\newcommand{\removequeries}{\renewcommand{\query}[1]{}}

%% web related requires: hyperref package
\newcommand{\webref}[2]{\href{#1}{\underline{#2}}}
\newcommand{\web}[1]{\webref{#1}{#1}}

%
% alphalist:  a), b), ...
%
\newcounter{alphanum}
\newenvironment{alphalist}{\begin{list}{\alph{alphanum})}
        {\usecounter{alphanum}\setlength{\leftmargin}{1.5em}}
        \rm}{\end{list}}

%% Grammar specific commands
\newcommand{\emptystring}{{\Large \tsym{\epsilon}}~}
\newcommand{\production}{\item}
\newcommand{\arr}{$\rightarrow$~}    %% arrow
\newcommand{\key}[1]{{\sffamily\bfseries #1}}  %% keyword
\newcommand{\tok}[1]{{\scshape\bfseries #1}}  %% token returned such as number
\newcommand{\non}[1]{{\it #1}}       %% nonterm
\newcommand{\sep}{~$\mid$~~}     %% alternation
\newcommand{\tsym}[1]{$\boldsymbol{#1}$}  %% token returned such as number

%% Walsh and Bit specific
%\newcommand{\contreq}[3]{#3 = {\cal C} (#1, #2)}
%\newcommand{\contr}[2]{{\cal C} (#1, #2)}
%\newcommand{\widevec}[1]{\overrightarrow{#1}}  too big and ugly
%\newcommand{\wt}{T}
\DeclareMathOperator*{\acos}{acos}
\DeclareMathOperator*{\asin}{asin}
\DeclareMathOperator*{\atan}{atan}
\DeclareMathOperator*{\argmax}{\arg\!\max}
\DeclareMathOperator*{\argmin}{\arg\!\min}
\DeclareMathOperator*{\bc}{bc}
\DeclareMathOperator*{\degray}{degray}
\DeclareMathOperator*{\deungray}{deungray}
\DeclareMathOperator*{\dimof}{dim}  % \dim is already defined
\DeclareMathOperator*{\extract}{extract}
\DeclareMathOperator*{\fitness}{fitness}
\DeclareMathOperator*{\gray}{gray}
\DeclareMathOperator*{\maxdim}{maxdim}
\DeclareMathOperator*{\pack}{pack}
\DeclareMathOperator*{\parity}{parity}
\DeclareMathOperator*{\pred}{pred}
\DeclareMathOperator*{\ungray}{ungray}
\DeclareMathOperator*{\unpack}{unpack}
\newcommand\degrees[1]{\ensuremath{#1^\circ}}
\newcommand{\betterthan}{\succ}
\newcommand{\binary}[1]{\mbox{\tt #1$_2$}}
\newcommand{\bits}[1]{\mbox{\tt #1}}
\newcommand{\bspace}{{\cal B}}
\newcommand{\chisq}{\mbox{$\chi^2$}}
\newcommand{\contains}{\supseteq}
\newcommand{\coversnot}{\not\supseteq}
\newcommand{\covers}{\supseteq}
\newcommand{\domain}{{\cal D}}
\newcommand{\dspace}{{\cal D}}
\newcommand{\embed}[2]{{\cal E} (#1, #2)}
\newcommand{\fhat}{\widehat{f}}
\newcommand{\forallsp}{\tnygap\forall~~}
\newcommand{\func}[1]{\mbox{\sl #1}}
\newcommand{\ghat}{\widehat{g}}
\newcommand{\given}{\,|\,}
\newcommand{\hhat}{\widehat{h}}
\newcommand{\isin}{\subseteq}
\newcommand{\isnotin}{\not\subseteq}
\newcommand{\logand}{\wedge}
\newcommand{\logiff}{\leftrightarrow}
\newcommand{\logimp}{\rightarrow}
\newcommand{\loglshift}{\stackrel{\mbox{\tiny shift}}{\longleftarrow}}
\newcommand{\logminus}{-}
\newcommand{\logeq}{=}
\newcommand{\lognand}{\uparrow}
\newcommand{\lognor}{\downarrow}
\newcommand{\logneg}{\sim\!\!}
\newcommand{\lognot}[1]{\overline{#1}}
\newcommand{\logor}{\vee}
\newcommand{\logpk}{\vdash}
\newcommand{\logrshift}{\stackrel{\mbox{\tiny shift}}{\longrightarrow}}
\newcommand{\logunpk}{\dashv}
\newcommand{\logxor}{\oplus}
\newcommand{\ones}{\vec{\tt 1}}
\newcommand{\pbits}[1]{{{\tt #1}'s}}
% \newcommand{\smallbc}{\mbox{\scriptsize\sl bc}}
% \newcommand{\smallpack}{\mbox{\scriptsize\sl pack}}
% \newcommand{\smallparity}{\mbox{\scriptsize\sl parity}}
\newcommand{\smallpm}{\raisebox{.3ex}{$\scriptstyle \pm$}}
% \newcommand{\smallunpack}{\mbox{\scriptsize\sl unpack}}
% \newcommand{\smbc}{\mbox{\scriptsize\sl bc}}
% \newcommand{\smpack}{\mbox{\scriptsize\sl pack}}
% \newcommand{\smparity}{\mbox{\scriptsize\sl parity}}
% \newcommand{\smpm}{\raisebox{.3ex}{$\scriptstyle \pm$}}
% \newcommand{\smunpack}{\mbox{\scriptsize\sl unpack}}
\newcommand{\strictcontains}{\supset}
\newcommand{\trm}[2]{\walsh_{#2}(#1)w_{#2}}
\newcommand{\unembed}[2]{{\cal C} (#1, #2)}
\newcommand{\vose}{\phi}
\newcommand{\walsh}{\psi}
\newcommand{\wcsum}{\wc_{sum}}
\newcommand{\wc}{{\cal K}}
\newcommand{\worsethan}{\prec}
\newcommand{\wt}{\mathcal{W}}
\newcommand{\yspace}{{\cal Y}}
\newcommand{\zeros}{\vec{\tt 0}}

\newcommand{\wfrac}{\frac{1}{2^L}}
\newcommand{\sumhp}[3]{\sum_{{#1}\elemof h_{#2, #3}}}
\newcommand{\sumbitlen}[1]{\sum_{{#1}=0}^{L-1}}
\newcommand{\sumlen}[1]{\sum_{{#1}=0}^{2^L-1}}
\newcommand{\sumsize}[1]{\sum_{{#1}=0}^{2^L-1}}  % don't use this anymore
\newcommand{\prodlen}[1]{\prod_{{#1}=0}^{L-1}}
\newcommand{\invpow}[1]{\inv{2^#1}}

%% Math
% \newcommand{\openm}{\vspace*{-.06in}\begin{equation}}
% \newcommand{\closem}{\vspace*{-.04in}\end{equation}}
% \newcommand{\openm}{$\displaystyle}
% \newcommand{\closem}{$}
\renewcommand{\emptyset}{\varnothing}
\newcommand{\onetoone}{$1\!\!-\!1$}
\newcommand{\reals}{\mathbb{R}}
\newcommand{\integers}{\mathbb{I}}
\newcommand{\posintegers}{\mathbb{Z}}
\newcommand{\modintegers}[1]{{{\posintegers}_{#1}}}

% logical operators in text form
\newcommand{\ortxt}{{\sc or}}
\newcommand{\andtxt}{{\sc and}}
\newcommand{\nottxt}{{\sc not}}
\newcommand{\xortxt}{{\sc exclusive-or}}
\newcommand{\Xortxt}{{\sc Exclusive-or}}
\newcommand{\truetxt}{{\sc true}}
\newcommand{\falsetxt}{{\sc false}}

\newcommand{\NP}{\mbox{NP}}
\newcommand{\infinity}{\infty}

\newcommand{\avg}[1]{\langle #1 \rangle}
\newcommand{\ceil}[1]{{\lceil #1 \rceil}}
% \newcommand{\comb}[2]{{\binomal{#1}{#2}}}
\newcommand{\comb}[2]{{#1\choose#2}}
\newcommand{\conv}{\circ}
\newcommand{\cross}{\!\times\!}
\newcommand{\dotprod}{\cdot}
\newcommand{\ds}{\displaystyle}
\newcommand{\elemof}{\in\,}
\newcommand{\notelemof}{\not\in\,}
\newcommand{\floor}[1]{{\lfloor #1 \rfloor}}
\newcommand{\intersect}{\cap}
\newcommand{\bigintersect}{\bigcap}
\newcommand{\inv}[1]{\frac{1}{#1}}
\newcommand{\mapto}{{\rightarrow}\,}
% \newcommand{\mod}{\bmod}
\newcommand{\order}[1]{{\cal{O}}(#1)}
\newcommand{\rounddown}[2]{{\lfloor #1 \rfloor \raisebox{-1.0ex}{\scriptsize $#2$}}}
\newcommand{\rspace}{{\cal R}}
\newcommand{\achspace}{{\cal H}}
\newcommand{\smtxt}[1]{{\scriptstyle\rm #1}}
\newcommand{\st}{\,:\,}
\newcommand{\sttxt}{~~~s.t.~~~}
\newcommand{\subsetof}{\subseteq}
\newcommand{\suchthat}{~:~}
\newcommand{\supermath}[1]{\raisebox{2ex}{$\scriptscriptstyle #1$}}
% \newcommand{\therefore}{.\!\!\cdot\!\!.}
\newcommand{\transpose}[1]{{#1}^{\scriptscriptstyle \sf T}}
\newcommand{\transposep}[1]{{(#1)}^{\scriptscriptstyle \sf T}}
\newcommand{\sminv}{\raisebox{.5ex}{\Large $-1$}}
\newcommand{\smtranspose}{\raisebox{.5ex}{\large\sf T}}
\newcommand{\txt}[1]{\mbox{~~#1~~}}
\newcommand{\tx}[1]{\mbox{#1}}
\newcommand{\ul}[1]{\underline{#1}}
\newcommand{\union}{\cup}
\newcommand{\bigunion}{\bigcup}

%% Sections
%Theorems
\newcounter{theoremnum}
\newcommand{\theorem}[1]{{\refstepcounter{theoremnum}\vspace*{.05in}\noindent\bfseries \underline{Theorem} \thetheoremnum\ } (#1) \index{#1 Theorem}\index{Theorems!#1}
\itshape
}

% theorem proving support
\newcommand{\case}[1]{\vspace*{.05in}\noindent{\bfseries\upshape CASE\ #1:}}
\newcommand{\numcase}[1]{\vspace*{.05in}\noindent{\bfseries\upshape CASE\ #1:}}
\newcommand{\numexample}[1]{\vspace*{.05in}\noindent{\bfseries\upshape EXAMPLE\ #1:}}
\newcommand{\derivation}{\vspace*{.05in}\noindent{\bfseries\upshape Derivation: }}
\newcommand{\theoremend}{\upshape}
% \newcommand{\proof}{\upshape \vspace*{-.05in}\noindent{\bfseries Proof: }}
\renewcommand{\qedsymbol}{\upshape \hspace*{0.9\columnwidth} $\Box$\hspace*{-1ex}\raisebox{.5ex}{$\Box$}}
\newcommand{\diamondqed}{\hspace*{0.9\columnwidth} $\Diamond$}
\newcommand{\blackqed}{$\blacksquare$}

%Lemmas
\newcounter{lemmanum}
\newcommand{\lemma}{{\refstepcounter{lemmanum}\vspace*{.05in}\noindent\bfseries \underline{Lemma} \thelemmanum\ }
\itshape}

% Conjectures
\newcounter{conjecturenum}
\newcommand{\conjecturenonum}{{\refstepcounter{conjecturenum}\vspace*{.05in}\noindent\bfseries \underline{Conjecture}\ }}
\newcommand{\conjecture}{{\refstepcounter{conjecturenum}\vspace*{.05in}\noindent\bfseries \underline{Conjecture} \theconjecturenum\ }
\itshape}

% Corollaries
\newcounter{corollarynum}[theoremnum]
\newcommand{\corollarynonum}{{\refstepcounter{corollarynum}\vspace*{.05in}\noindent\bfseries \underline{Corollary}\ }}
\newcommand{\corollary}{\refstepcounter{corollarynum}\vspace*{.05in}\noindent{\bfseries \underline{Corollary} \thetheoremnum \alph{corollarynum}\ }
\itshape}
\newcommand{\lastcorollarynum}{\thetheoremnum \alph{corollarynum}}

\newcommand{\namedcorollary}[1]{{\refstepcounter{corollarynum}\vspace*{.05in}\noindent\bfseries \underline{Corollary} \thetheoremnum \alph{corollarynum}\ } (#1) \index{#1 Corollary}\index{Corollarys!#1}}


% Observations
\newcommand{\observations}{\noindent{\bfseries Observations: }}

% Examples
\newenvironment{example}{\begin{quote}{\bf Example:~}}{\end{quote}}

%% General
\newcommand{\note}[1]{{\large\bfseries Note: #1}}
%\newcommand{\note}[1]{\fbox{\ #1\ }}
\newcommand{\abeam}[1]{\rule{#1}{0.0in}}
\newcommand{\citehere}{\fbox{\ add cite here\ }}
\newcommand{\defit}[1]{{\bfseries #1\/}\index{#1}}
\newcommand{\downstrut}[1]{\rule[-#1]{0pt}{#1}}
\newcommand{\upstrut}[1]{\rule{0.0in}{#1}}
\newcommand{\hr}{\noindent\rule{\columnwidth}{0.1\baselineskip}\\}
\newcommand{\indexit}[1]{{#1}\index{#1}}
\newcommand{\numnd}[1]{{$ #1^{nd} $}}
\newcommand{\numrd}[1]{{$ #1^{rd} $}}
\newcommand{\numst}[1]{{$ #1^{st} $}}
\newcommand{\numth}[1]{{$ #1^\textrm{th} $}}

\newcommand{\neggap}{\hspace*{-2ex}}
\newcommand{\exgap}{\hspace*{1ex}}
\newcommand{\tnygap}{\hspace*{.125in}}
\newcommand{\smgap}{\hspace*{.25in}}
\newcommand{\gap}{\hspace*{.5in}}
\newcommand{\lggap}{\hspace*{1.0in}}
\newcommand{\hugegap}{\hspace*{2.0in}}

\newcommand{\negvgap}{\vspace*{-2ex}}
\newcommand{\tnyvgap}{\vspace*{1ex}}
\newcommand{\smvgap}{\vspace*{2ex}}
\newcommand{\vgap}{\vspace*{4ex}}
\newcommand{\lgvgap}{\vspace*{8ex}}
\newcommand{\hugevgap}{\vspace*{16ex}}

%% Cheap Slide making
%\newcommand{\graphic}[2]{\fbox{\rule{0.0in}{#1}\large\rm {\bfseries #2 }}}
%\newcommand{\graphic}[2]{\framebox[0.99\textwidth]{\rule{0.0in}{#1}\large\rm {\bfseries #2 }}}
%\newcommand{\graphic}[2]{{\rule{0.0in}{#1}\large\rm {\bfseries #2 }}}
%\newcommand{\help}[1]{{\large\rm {\bfseries [} #1{\bfseries ]} }}



\newcommand{\graphic}[2]{\makebox[0.99\textwidth]{\rule{0.0in}{#1}\large\rm {\bfseries #2 }}}
\newcommand{\defhead}[1]{{\huge\sf \underline{#1}~}}
\newcommand{\newslide}{\clearpage\Huge\sf}
\newcommand{\slidehead}[1]{\clearpage{\Huge\sf \underline{#1}{\vspace*{.2in}}}}
\newcommand{\slideheadextra}[1]{{\vspace*{-.2in}}{\Huge\sf \underline{#1}}}
\newcommand{\slidesub}[1]{{\huge\sf #1{\vspace*{.1in}}}}
\newcommand{\say}[1]{{}}
\newcommand{\slidenote}[1]{}
\newcommand{\smbits}[1]{\mbox{\large\tt #1}}
\newcommand{\subhead}[1]{{\vspace*{.1in}\Huge\sf #1{\vspace*{.05in}}}}
\newcommand{\titlepagebox}[1]{\clearpage{\Huge\bfseries\sf
\vspace*{.38\textheight}
\setlength{\fboxrule}{2pt}\setlength{\fboxsep}{5pt}
\centerline{\fbox{\usefont{OT1}{lcmss}{m}{n}\fontsize{30pt}{36pt}\selectfont #1}}}}

% example of common parameters
% \textwidth=7.5in
% \textheight=9in
% \columnsep=.25in
% \oddsidemargin=0in
% \topmargin=-.5in
% \pagestyle{empty}
% \pagestyle{plain}


\parskip=.9ex
\textwidth=7.0in
\textheight=9.0in
\oddsidemargin=-.25in
%\oddsidemargin=0in
\topmargin=-.75in

\renewcommand{\arraystretch}{1.2}
\begin{document}

\title{SLang}
\author{Robert B. Heckendorn
  \and
Damian Ball 
  \and 
updated by Keith Drew}
\maketitle

The command language for the simulator is called \textbf{slang}.  Here is the
syntax and semantics (recovered from Damian's thesis, updated by Keith):

\begin{enumerate}
\production \non{program} \arr \non{command-list}

\production \non{command-list} \arr \non{command-list} \non{command} \sep \non{command}

\production \non{command} \arr \non{exit} \sep \non{hi} \sep \non{city} \sep \non{edge} \sep \non{agent} \sep  \non{rmedge} \sep \non{run} \sep \non{show} \sep \non{safety} \sep \non{node}

\production \non{exit} \arr \key{exit} \sep \key{quit}

Quits the program.

\production \non{hi} \arr \key{hi}

Prints ``Howdy!'' as a test that the program is ``awake'' and ``listening'' to command input.

\production \non{city} \arr \key{city} \tok{STRING} \tok{INT} \tok{INT}

Creates a city object with:
\begin{itemize}
\item the title given by the string 
\item fixed number, $N$, of nodes numbered $1$ to $N$
\item maximum out degree for any of the nodes so pointers the edges adjacent to the node can be
allocated once and have to manage variable length lists in this speed critical function.
\end{itemize}

\production \non{edge} \arr \key{edge} \tok{INT} \tok{INT} \tok{FLOAT} \tok{FLOAT} \tok{FLOAT} \tok{FLOAT} \tok{FLOAT}\\


This adds an edge in the current city.   The edge is identified by (from node, to node, order of entry).  The time spent by a car on an edge is determined by the
parameters given for each edge.   The time is:
\[
\textrm{time} = \mathcal{F} \left( 1 + b \left(\frac{c}{c_{max}}\right)^\beta \right)
\]
where $c$ is the amount of traffic on the edge, $c_{max}$ is the capacity,
$\mathcal{F}$ is the free flow time, b and $\beta$ are constants.
The parameters in order are the following.   If parameters are not given it defaults to the last used value for that positional parameter (not true for updated version - parameters must be present in SLANG file).
\begin{itemize}
\item From node number
\item To node number
\item $c_{max}$: Capacity as the number of cars per unit time for the road.  Used in the link cost formulas.
\item $\mathcal{F}$: Free flow travel time, the amount of unit time that it takes a car to travel the length of the edge under ideal conditions
\item $b$: coefficient in the link cost formulas, also commonly named $\alpha$
\item $\beta$: power or exponent in the link cost formulas
\item The last value specified is the safety value of the specified edge, typically ranging from $0-1$.

%\item Length of the road in whatever units you want.  \textbf{Not currently} needed for our optimizations.
%\item Speed limit.  \textbf{Not currently} needed for our optimizations.
%\item Cost to travel a link.  \textbf{Not currently} needed for our optimizations.
% \item Type of link/zone.  \textbf{Not currently} needed for our optimizations.
\end{itemize}

\production \non{rmedge} \arr \key{rmedge} \tok{INT} \tok{INT} \sep \\
\smgap  \key{rmedge} \tok{INT} \tok{INT} \tok{INT}

Deletes the edge identified by the from and to node numbers or the from and to node numbers plus
the order in which they were added.

\production \non{node} \arr \key{node} \tok{INT} \tok{INT} \tok{FLOAT} \tok{FLOAT} \tok{FLOAT} \tok{FLOAT}

Specifies node parameters for the node indicated with the following parameters:
\begin{itemize}
	\item Node ID
	\item Node capacity
	\item Agent wait time for node - the time that an agent will wait at the node before being re-evaluated (in unit time)
	\item The safety of the given node
	\item The node longitude
	\item The node latitude
\end{itemize}

\production \non{agent} \arr \key{agent} \tok{INT} \tok{INT} \tok{INT}

Specifies simulation agents with the following parameters:
\begin{itemize}
	\item The starting node of the agent in each simulation (evaluation).
	\item The number of agents to start at the specified node.
	\item The number of individuals in the agent group.
\end{itemize}

\production \non{probabilities} \arr \key{probabilities} \tok{STRING}

STRING specifies a file to read for a population of probability sets that should be used to initialize the population before running a new evolution.

\production \non{anyint} \arr \tok{INT} \sep \key{*}

%\production \non{od} \arr \key{od} \non{anyint} \non{anyint}  \tok{FLOAT} \tok{FLOAT}


%\begin{itemize}
%\item The first two arguments are the From and To nodes.  They could also be \key{*}'s which
%indicate for all nodes in that position.   e.g.  od * 17 3.14 2.78 means for all nodes going to
%node 17 the values of mean and standard deviation are 3.14 and 2.78 respectively.
%\item Mean number of cars per hour generated going from the From node to the To node.
%The mean is really a proportion.   
%\item Standard Deviation of the cars per hour.
%\end{itemize}
%The default value is no cars are generated going between any pair of nodes.
%Any values assigned for cars going from a node to itself will not be used in generating car schedules.
%Values of already defined OD pairs may be changed by just specifying the same OD pair 
%and the new values.
%For instance all the OD pairs could be set to a base value and then some OD pairs of interest
%could be set to zero or to a higher value afterward.


\production \non{run} \arr \key{run} \tok{INT} \tok{INT} \tok{INT}

Initialize the simulation and run the simulation and learning.   The parameters are:
\begin{itemize}
\item Total number of evolution generations.
\item The number of time steps between reporting progress. Must evenly divide simulation steps.
\item The number of runs of the whole optimization, restarting each time.
\end{itemize}

\production \non{simulate} \arr \key{simulate} \tok{STRING}

STRING specifies a file with probabilities that should be run through the traffic simulation. This is used to evaluate fitness of a human-written set of probabilities. Use for testing simple cases to prove the fitness function works, and to prove the model can solve simple problems. 

\production \non{slang} \arr \key{slang} \tok{STRING}

Command to output a SLang file based on the current city topology. STRING designates the filename to write the SLang file to.

%\production \sout{\non{go} \arr \key{go} \tok{INT} \tok{INT} \sep \key{go} \tok{INT} \sep \key{go}}

%Run the simulation using the default parameters from the last \key{run} or \key{go} command.

%\production \sout{\non{eval} \arr \key{eval} TBD...}

%Evaluate the performance of cars running the learned pheromone tables.
%This runs the simulation collecting statistics on
%performance without altering the solution to the optimization.

%\production \non{show} \arr \non{city} \sep \non{edges} \sep \non{topology} \sep \non{od} \sep \non{pheromones} \sep \non{dotformat} \sep \non{edgeload} \sep \non{heatmap} \sep \non{allpaths}
\production \non{show} \arr \non{edges} \sep \non{topology} \sep \non{probabilities} \sep \non{drawable} \sep \non{parameters}

Output commands show edges, city topology, and probabilities. The drawable command outputs the city topology information needed to create a visual representation of the graph. Parameters refer to the evolution parameters (generations, sigma, population size, etc)

%\production \non{show city} \arr \key{show city}

\production \non{show edges} \arr \key{show edges}

\production \non{show topology} \arr \key{show topology}

%\production \non{show od} \arr \key{show od}

%\production \non{show pheromones} \arr \key{show pheromones} \tok{INT}
\production \non{show probabilities} \arr \key{show probabilities}

\production \non{show drawable} \arr \key{show drawable}

%\production \non{show dotformat} \arr \key{show dot} \tok{INT}

%\production \non{show edgeload} \arr \key{show edgeload}

%\production \non{show heatmap} \arr \key{show heatmap}

%\production \non{show allpaths} \arr \key{show allpaths}
\bigskip
\bigskip

\textbf{\LARGE{NOTE:}} The grammar rules here indicate the theoretical requirements of the grammar. Rules may need to change depending on implementation and optimization algorithm. SLang was created with Ant-Colony Optimization in mind, and has been repurposed here for use with an Evolution Strategy used to optimize a cloud of traffic-assignment probabilities. 

\end{enumerate}
\end{document}
